\documentclass[10pt]{article}
\usepackage[utf8]{inputenc}
\usepackage[T1]{fontenc}
\usepackage{amsmath}
\usepackage{amsfonts}
\usepackage{amssymb}
\usepackage[version=4]{mhchem}
\usepackage{stmaryrd}
\usepackage{bbold}
\usepackage{mathrsfs}

\begin{document}

\begin{align*}
\mathbf{o}(h) & \geq \mathbb{E}[X(t+h)-X(t) \mid X(t)=x]  \tag{1}\\
f(t, x) & =\lim _{h \rightarrow 0} \frac{1}{h} \mathbb{E}[X(t+h)-X(t) \mid X(t)=x]  \tag{2}\\
g^{2}(t, x) & =\lim _{h \rightarrow 0} \frac{1}{h} \mathbb{E}\left[(X(t+h)-X(t))^{2} \mid X(t)=x\right] \tag{3}
\end{align*}


Lemma 1. Assume $p(t, y \mid s, x) \in C^{1,2}((0, \infty) \times(0,1))$ and $X(t)$ satisfies above conditions. Then for any $(s, x) \in(0, \infty) \times(0,1)$ and $(t, y) \in(s, \infty) \times(0,1)$, the transition probability density $p(t, y \mid s, x)$ satisfies the Fokker-Planck equation

$$
\frac{\partial}{\partial t} p(t, y \mid s, x)=-\frac{\partial}{\partial y}(f(t, y) p(t, y \mid s, x))+\frac{1}{2} \frac{\partial^{2}}{\partial^{2} y}\left(g^{2}(t, y) p(t, y \mid s, x)\right)
$$

with the boundary condition $p(s, y \mid s, x)=\delta(y-x)$, where $\delta$ is the Dirac function.\\
Proof. Since $p(t, y \mid s, x)$ is the transition probability density function, $p(s, y \mid s, x)=$ $\delta(y-x)$ holds automatically. Next, for any $\phi \in \mathscr{C}_{c}^{2}(0,1)$ we estimate

$$
\begin{aligned}
\mathscr{I} & :=\int_{0}^{1} \frac{\partial}{\partial t} \phi(y) p(t, y \mid s, x) \mathrm{d} y=\frac{\partial}{\partial t} \int_{0}^{1} \phi(y) p(t, y \mid s, x) \mathrm{d} y \\
& =\lim _{h \rightarrow 0} \frac{1}{h}\left(\int_{0}^{1} \phi(y) p(t+h, y \mid s, x) \mathrm{d} y-\int_{0}^{1} \phi(z) p(t, z \mid s, x) \mathrm{d} z\right)
\end{aligned}
$$

To that end, first use the Kolmogorov-Chapman equation and Fubini's Theorem to obtain


\begin{align*}
\mathscr{I} & =\lim _{h \rightarrow 0} \frac{1}{h}\left(\int_{0}^{1} \phi(y)\left(\int_{0}^{1} p(t+h, y \mid t, z) p(t, z \mid s, x) d z\right) \mathrm{d} y-\int_{0}^{1} \phi(z) p(t, z \mid s, x) \mathrm{d} z\right) \\
& =\lim _{h \rightarrow 0} \frac{1}{h}\left(\int_{0}^{1}\left(\int_{0}^{1} \phi(y) p(t+h, y \mid t, z) \mathrm{d} y\right) p(t, z \mid s, x) \mathrm{d} z-\int_{0}^{1} \phi(z) p(t, z \mid s, x) \mathrm{d} z\right) \\
& =\lim _{h \rightarrow 0} \frac{1}{h} \int_{0}^{1}\left(\int_{0}^{1} \phi(y) p(t+h, y \mid t, z) \mathrm{d} y-\phi(z)\right) p(t, z \mid s, x) \mathrm{d} z \\
& =\int_{0}^{1}\left(\lim _{h \rightarrow 0} \frac{1}{h} \int_{0}^{1}(\phi(y)-\phi(z)) p(t+h, y \mid t, z) \mathrm{d} y\right) p(t, z \mid s, x) \mathrm{d} z \tag{4}
\end{align*}


To further estimate $\mathscr{I}$, we split the inner integral in (4) into two parts. More precisely, let $\varepsilon$ be an arbitrary positive number and write


\begin{align*}
\int_{0}^{1}(\phi(y)-\phi(z)) p(t+h, y \mid t, z) \mathrm{d} y= & \int_{|z-y|>\varepsilon}(\phi(y)-\phi(z)) p(t+h, y \mid t, z) \mathrm{d} y \\
& +\int_{|z-y| \leq \varepsilon}(\phi(y)-\phi(z)) p(t+h, y \mid t, z) \mathrm{d} y \tag{5}
\end{align*}


Then by using the property of the probability density of a diffusion process, there exists $K=K(\phi)$ such that


\begin{equation*}
\int_{|z-y|>\varepsilon}(\phi(y)-\phi(z)) p(t+h, y \mid t, z) \mathrm{d} y \leq K \int_{|z-y|>\varepsilon} p(t+h, y \mid t, z) \mathrm{d} y \leq \mathbf{o}(h) \tag{6}
\end{equation*}


On the other side, by Taylor's expansion of $\phi$ at $z$ and (2) - (3) again,


\begin{align*}
& \frac{1}{h} \int_{|z-y| \leq \varepsilon}(\phi(y)-\phi(z)) p(t+h, y \mid t, z) \mathrm{d} y \\
= & \frac{1}{h} \int_{|z-y| \leq \varepsilon}\left(\phi^{\prime}(z)(y-z)+\frac{1}{2} \phi^{\prime \prime}(z)(y-z)^{2}\right) p(t+h, y \mid t, z) \mathrm{d} y+\mathbf{o}\left(\varepsilon^{2}\right) \\
= & \phi^{\prime}(z) f(t, z)+\frac{1}{2} \phi^{\prime \prime}(z) g^{2}(t, z)+\mathbf{o}\left(\varepsilon^{2}\right) \tag{7}
\end{align*}


Applying estimates (6) and (7) to (5) and letting $\varepsilon \rightarrow 0$ gives


\begin{equation*}
\lim _{h \rightarrow 0} \frac{1}{h} \int_{0}^{1}(\phi(y)-\phi(z)) p(t+h, y \mid t, z) \mathrm{d} y=\phi^{\prime}(z) f(t, z)+\frac{1}{2} \phi^{\prime \prime}(z) g^{2}(t, z) \tag{8}
\end{equation*}


Now inserting (8) into (4) and integrating by parts results in

$$
\begin{aligned}
\int_{0}^{1} \frac{\partial}{\partial t} \phi(y) p(t, y \mid s, x) \mathrm{d} y & =\int_{0}^{1}\left(f(t, z) \phi^{\prime}(z)+\frac{1}{2} g^{2}(t, z) \phi^{\prime \prime}(z)\right) p(t, z \mid s, x) \mathrm{d} z \\
& =\int_{0}^{1}\left(-\frac{\partial}{\partial z}(a(t, z) p(t, z \mid s, x))+\frac{\partial^{2}}{\partial^{2} z}\left(\frac{1}{2} b^{2}(t, z) p(t, z \mid s, x)\right)\right) f(z) d z \\
& \left.=\int_{0}^{1}\left(-\frac{\partial}{\partial y}(f(t, y) p(t, y \mid s, x))+\frac{\partial^{2}}{\partial^{2} y}\left(\frac{1}{2} g^{2}(t, y) p(t, y \mid s, x)\right)\right) \phi(y) \mathrm{d} \rightsquigarrow 10\right)
\end{aligned}
$$

Since $\phi$ is an arbitrary function in $C_{c}^{2}(0,1)$ and $p(t, y \mid s, x)$ is continuous, (10) implies that


\begin{equation*}
\frac{\partial}{\partial t} p(t, y \mid s, x)=-\frac{\partial}{\partial y}(f(t, y) p(t, y \mid s, x))+\frac{1}{2} \frac{\partial^{2}}{\partial^{2} y}\left(g^{2}(t, y) p(t, y \mid s, x)\right) \tag{11}
\end{equation*}


for all $(s, x) \in(0, \infty) \times(0,1)$ and $(t, y) \in(s, \infty) \times(0,1)$. The proof is complete.


\end{document}